\documentclass[11pt,a4paper]{amsart}
\usepackage{amsmath,amssymb,amsthm,mathtools}
\usepackage[margin=1.2in]{geometry}
\usepackage{hyperref}

\theoremstyle{plain}
\newtheorem{theorem}{Theorem}[section]
\newtheorem{lemma}[theorem]{Lemma}
\theoremstyle{definition}
\newtheorem{definition}[theorem]{Definition}
\theoremstyle{remark}
\newtheorem{remark}[theorem]{Remark}
\newtheorem{conjecture}[theorem]{Conjecture}

\newcommand{\lean}[1]{\hfill\texttt{\footnotesize #1}}
\newcommand{\ZZ}{\mathbb{Z}}
\newcommand{\QQ}{\mathbb{Q}}
\newcommand{\NN}{\mathbb{N}}
\newcommand{\vp}{v_p}
\DeclareMathOperator{\lcm}{lcm}

\title{A closed form and a Lucas congruence for Ap\'ery's companion sequence}
\author{[Author]}
\date{}

\begin{document}
\begin{abstract}
We give a two-term closed form for the numerators $b_n$ of Ap\'ery's approximations
to $\zeta(3)$, a Lucas-type congruence $p^3b_{ap+r}\equiv b_a\,a_r\pmod p$
for them, and a correction to a published integrality claim. Every theorem and lemma
below is proved in Lean~4 and checked by its kernel, and is followed by the name of the
declaration that proves it; the one conjecture is labelled as such.
\end{abstract}
\maketitle

\section{Introduction}

Ap\'ery's proof that $\zeta(3)\notin\QQ$ rests on the pair of sequences
$a_n=\sum_k\binom nk^2\binom{n+k}k^2$ and a companion $b_n$ with $b_n/a_n\to\zeta(3)$.
The first sequence is arithmetically well understood: it satisfies a Lucas congruence
$a_{ap+r}\equiv a_aa_r\pmod p$ (Gessel \cite{Gessel1982}), so it respects base-$p$ digits.
The companion $b_n$ has attracted no such statement, and it is usually written with a
nested alternating inner sum, which makes it awkward to handle modulo $p$.

We report three things.

First, a closed form (Theorem \ref{thm:closed}): $b_n$ is the same binomial sum as $a_n$
weighted by a \emph{two-term} harmonic expression,
\[
  b_n=\sum_{k=0}^n\binom nk^2\binom{n+k}k^2\bigl(2H^{(3)}_n-H^{(3)}_k\bigr),
  \qquad H^{(r)}_m=\sum_{j=1}^m j^{-r}.
\]
The nested sum disappears.

Second, a Lucas-type congruence for the companion itself (Theorem \ref{thm:bA}):
\[
  p^3\,b_{ap+r}\equiv b_a\,a_r \pmod p \qquad (a,r<p).
\]
The factor $p^3$ is the weight of $\zeta(3)$. The same argument gives the analogous
congruence for Franel's companion at weight $2$ (Theorem \ref{thm:franel}).

Third, a correction. Brown and Zudilin \cite{BZ} record $d_n^{\,5}P_n\in\ZZ$ for the
numerators of their cellular approximations to $\zeta(5)$, where $d_n=\lcm(1,\dots,n)$.
Their own printed value $P_2=1190161/384$ does not satisfy this. At $n=2$ the defect is
exactly $12=2^2\cdot3$ (Theorem \ref{thm:twelve}); consequently any constant $c$ for
which $c\,d_n^{\,5}P_n$ is always integral is divisible by $12$. Whether $12$ suffices
for every $n$ we state only as Conjecture \ref{conj:twelve}: it is supported by
computation and by a proof for the primes $p\ge5$ given elsewhere, and it is not
formalised.

Finally we record a transcription slip in the same author's earlier work
(Section \ref{sec:errata}); it is benign, and we do not formalise it.

Throughout, $p$ is prime, and for $x,y\in\QQ$ we write $x\equiv y\pmod p$ to mean
$\vp(x-y)\ge1$.
\lean{PInt, PDvd, PCong}

\section{Ap\'ery's companion}

\begin{definition}
$\displaystyle A(n,k)=\binom nk^2\binom{n+k}k^2$ \ and \ $a_n=\sum_{k=0}^n A(n,k)$.
\lean{A, apery}
\end{definition}

\begin{lemma}\label{lem:Adigits}
For $r,s<p$, \ $A(ap+r,\,bp+s)\equiv\binom ab^2\binom{a+b}b^2\,A(r,s)\pmod p$.
\lean{A\_digits}
\end{lemma}

This is Lucas' theorem applied to each binomial, together with the fact that a carry
$r+s\ge p$ makes both sides vanish.
\lean{choose\_digits, choose\_carry\_zero}

\begin{definition}\label{def:bApery}
$P(n)=34n^3+51n^2+27n+5$, and $b_n$ is defined by
\[
  b_0=0,\qquad b_1=6,\qquad (n+2)^3b_{n+2}=P(n+1)\,b_{n+1}-(n+1)^3b_n .
\]
\lean{Ppol, bApery}
\end{definition}

This is Ap\'ery's recurrence; $b_2=351/4$ and $b_3=62531/36$.

\begin{definition}
$\displaystyle \tilde b_n=\sum_{k=0}^n A(n,k)\bigl(2H^{(3)}_n-H^{(3)}_k\bigr)$.
\lean{Harm, bMin}
\end{definition}

\begin{lemma}\label{lem:rec}
$(m+2)^3\tilde b_{m+2}-P(m+1)\,\tilde b_{m+1}+(m+1)^3\tilde b_m=0$ for all $m\ge0$.
\lean{bMin\_rec}
\end{lemma}

The proof is a single telescoping step: applying the recurrence operator to the summand
of $\tilde b$ produces an exact difference $\Psi(m,k+1)-\Psi(m,k)$, whose sum collapses to
the boundary values $\Psi(m,0)=\Psi(m,m+3)=0$. The two Zeilberger-type certificates
involved are polynomial identities in $\QQ[m,k]$.
\lean{star, star\_pre}

\begin{theorem}\label{thm:closed}
$\tilde b_n=b_n$ for all $n$; that is,
\[
  \sum_{k=0}^n\binom nk^2\binom{n+k}k^2\bigl(2H^{(3)}_n-H^{(3)}_k\bigr)=b_n .
\]
\lean{bMin\_eq\_bApery, apery\_b\_harmonic\_closed\_form}
\end{theorem}

\begin{theorem}\label{thm:bA}
For $a,r<p$, \ $p^3\,b_{ap+r}\equiv b_a\,a_r\pmod p$.
\lean{bApery\_lucas}
\end{theorem}

The weight $3$ enters as follows: modulo $p$ the summands with $k\not\equiv0$ contribute
nothing after multiplication by $p^3$, while the layer $k=jp$ reproduces
$H^{(3)}$ at $\lfloor\,\cdot\,/p\rfloor$, since $p^3\sum_{m\le y,\ p\mid m}m^{-3}
=\sum_{j\le y/p}j^{-3}$. The remaining digit factorisation is Lemma \ref{lem:Adigits}.

\begin{theorem}\label{thm:franel}
Let $f_n=\sum_k\binom nk^3$ and
$\beta_n=\sum_k\binom nk^3\bigl(\tfrac14H^{(2)}_k+\tfrac34(H^{(1)}_k)^2
-\tfrac34H^{(1)}_kH^{(1)}_{n-k}\bigr)$. Then for $p\neq2$ and $a,r<p$,
\[
  p^2\,\beta_{ap+r}\equiv\beta_a\,f_r\pmod p .
\]
\lean{franel, bFranel, bFranel\_lucas}
\end{theorem}

\section{A denominator defect}

\begin{definition}
$d_n=\lcm(1,\dots,n)$, and $P_2=1190161/384$ is the value printed in \cite{BZ}.
\lean{dlcm, BZ\_P2}
\end{definition}

\begin{theorem}\label{thm:twelve}
$d_2^{\,5}P_2=1190161/12\notin\ZZ$, while $12\,d_2^{\,5}P_2=1190161\in\ZZ$; and if
$c\in\NN$ satisfies $c\,d_2^{\,5}P_2\in\ZZ$, then $12\mid c$.
\lean{BZ\_P2\_not\_integral, BZ\_P2\_factor\_twelve, BZ\_twelve\_minimal}
\end{theorem}

Here $d_2=2$ and $1190161\equiv1\pmod{12}$, so at $n=2$ the defect is exactly $2^2\cdot3$.
Specialising to $n=2$, any $c\in\NN$ with $c\,d_n^{\,5}P_n\in\ZZ$ for all $n$ is divisible
by $12$.

\begin{conjecture}\label{conj:twelve}
$12\,d_n^{\,5}P_n\in\ZZ$ for every $n\ge1$.
\end{conjecture}

Theorem \ref{thm:twelve} shows that no smaller constant can serve. Conjecture
\ref{conj:twelve} itself is \emph{not} formalised: it is supported by exact computation,
and its restriction to the primes $p\ge5$ is proved elsewhere, conditionally on two
finite identities verified but not machine-checked.

\section{An erratum}\label{sec:errata}

This section is expository and is not formalised.

Section~8 of \cite{Zu04} proves that one of $\zeta(5),\zeta(7),\zeta(9),\zeta(11)$ is
irrational. Its linear forms are built from a very-well-poised rational function, defined
at (8.2) by
\[
  \widetilde R(t)=(h_0+2t)\,
  \frac{\Gamma(h_0+t)^r\prod_{j=1}^{q}\Gamma(h_j+t)}
       {\Gamma(1+t)^r\prod_{j=1}^{q}\Gamma(1+h_0-h_j+t)},
\]
and then normalised arithmetically as $R=N\widetilde R$ with
$N=\prod_{j=r+1}^{q}(h_0-2h_j)!\big/\prod_{j=1}^{r}(h_j-1)!^2$. Equation~(8.7), which
restates $R(t)$ as an explicit product, omits the factor $(h_0+2t)$. The omission is in
the published text and in the arXiv source, not in a transcription of it.

The factor is needed. It is the unique factor of $\widetilde R$ that is antisymmetric
under $t\mapsto -t-h_0$, so it is what makes the well-poised symmetry (8.5),
\[
  \widetilde R(-t-h_0)=-\widetilde R(t),
\]
hold; and that symmetry is what forces $B_{jk}=(-1)^jB_{j,h_0-k}$ for the partial-fraction
coefficients, hence the vanishing of the coefficients of the \emph{even} zeta values in
$F(h)=\sum_j A_{j-1}\zeta(j-1)-A_0$, which is the content of Lemma~19 there.

The effect of the omission is not a loss of precision but a reversal of parity. Taking the
paper's own parameters $r=3$, $q=13$, $\eta=(91;27,27,27,29,30,\dots,38)$, $h_0=91n+2$,
$h_j=\eta_jn+1$, and computing the expansion of $F$ at $n=1$ in exact arithmetic: with
$(h_0+2t)$ present, the coefficients of $\zeta(3)$ and of $\zeta(4),\zeta(6),\zeta(8),
\zeta(10),\zeta(12)$ are all exactly zero and only $\zeta(5),\zeta(7),\zeta(9),\zeta(11)$
survive, as Theorem~3 of \cite{Zu04} requires; reading (8.7) literally, the even zeta
values survive instead and the odd ones vanish.

So (8.7) should read $R(t)=(h_0+2t)\cdot[\text{the printed product}]$. Equations (8.2),
(8.5) and (8.6) are consistent with the factor present, and the results of \S8 are
unaffected.

\section{The development}

Lean~4 (\texttt{v4.33.0-rc1}) with Mathlib \texttt{cd580e54}; ten modules, $2107$ lines;
no \texttt{sorry} and no \texttt{native\_decide}. Besides the results above it contains
Gessel's congruence $a_{ap+r}\equiv a_aa_r$ \cite{Gessel1982} and its analogue for the
Brown--Zudilin double sum
\[
  Q_n=\sum_{k,l}\binom{n+k}n\binom nk^2\binom{n+l}n\binom nl^2\binom{n+k+l}n,
\]
in both cases with the base-$p$ digit-product form
(\texttt{apery\_lucas}, \texttt{Q\_lucas}). Every declaration cited above satisfies
\[
  \texttt{\#print axioms}\;\longrightarrow\;
  [\texttt{propext},\ \texttt{Classical.choice},\ \texttt{Quot.sound}].
\]
Build with \texttt{lake exe cache get \&\& lake build}.

\begin{thebibliography}{9}
\bibitem{BZ}
F.~Brown and W.~Zudilin, \emph{On cellular rational approximations to $\zeta(5)$},
preprint \texttt{arXiv:2210.03391v3} (2022).
\bibitem{Gessel1982}
I.~M.~Gessel, \emph{Some congruences for Ap\'ery numbers},
J.\ Number Theory \textbf{14} (1982), no.~3, 362--368.
\bibitem{mathlib}
The mathlib Community, \emph{The Lean mathematical library}, CPP 2020, 367--381.
\texttt{arXiv:1910.09336}.
\bibitem{lean4}
L.~de Moura and S.~Ullrich, \emph{The Lean 4 theorem prover and programming language},
CADE~28, LNCS 12699, 625--635.
\bibitem{Zu04}
W.~Zudilin, \emph{Arithmetic of linear forms involving odd zeta values},
J.\ Th\'eor.\ Nombres Bordeaux \textbf{16} (2004), 251--291.
\end{thebibliography}

\end{document}
